\item \textbf{Clasificación de estrellas por luminosidad}.

\begin{tcolorbox}[colback=blue!2!white,colframe=blue!55!black,title={Enunciado}]
La luminosidad relativa de una estrella de masa \(M\) (en masas solares) se estima mediante la relación:
\[
L = M^{3.5}
\]

\begin{itemize}
    \item Escriba un programa que:
    \begin{itemize}
        \item Solicite al usuario la masa \(M\) de una estrella.
        \item Implemente una función \texttt{luminosidad(M)} que calcule \(L\).
        \item Utilice un ciclo \texttt{while} para procesar múltiples valores de masa hasta que el usuario ingrese \(-1\).
        \item Utilice una estructura \texttt{if/else} para indicar si la estrella es más o menos luminosa que el Sol.
    \end{itemize}
\end{itemize}
\end{tcolorbox}

\begin{solution}

\subsubsection*{Idea}
El problema consiste en calcular la luminosidad de una estrella a partir de su masa usando la relación \(L = M^{3.5}\). 

Se requiere:
\begin{itemize}
    \item Definir una función que reciba la masa \(M\) y retorne la luminosidad \(L\).
    \item Pedir valores al usuario repetidamente usando un ciclo \texttt{while}.
    \item Terminar el ciclo cuando el usuario ingrese \(-1\).
    \item Comparar la luminosidad obtenida con la del Sol (\(L=1\)) para clasificar la estrella.
\end{itemize}

\subsubsection*{Implementación}

\begin{pythonjupyter}
# Función que calcula la luminosidad
def luminosidad(M):
    return M**3.5

# Bucle principal
while True:
    # Solicitar masa
    M = float(input("Ingrese la masa de la estrella (-1 para salir): "))
    
    # Condición de término
    if M == -1:
        print("Programa finalizado.")
        break
    
    # Calcular luminosidad
    L = luminosidad(M)
    
    # Clasificación
    if L > 1:
        print(f"Luminosidad = {L:.2f} → Más luminosa que el Sol")
    else:
        print(f"Luminosidad = {L:.2f} → Menos luminosa que el Sol")
\end{pythonjupyter}

\subsubsection*{Resultado}
El programa permite ingresar múltiples masas de estrellas, calcular su luminosidad y clasificarlas en relación con el Sol, utilizando funciones, ciclos y estructuras condicionales.

\end{solution}
